% СПО ЛКС - вводная лекция
% Пынькин Д.А. (с) 2011
\documentclass[ignorenonframetext, hyperref={pdftex, unicode}]{beamer}
\usepackage{beamerthemesplit}

\usetheme{Pittsburgh}
\usecolortheme{dolphin}

\usepackage[russian]{babel}
\usepackage[utf8]{inputenc}
\usepackage[T1]{fontenc}
\usepackage{ulem}

\title[СПОЛКС]{Системое программное обеспечение локальных компьютерных сетей}
\subtitle{Введение}
\author{Денис Пынькин}
\date{2011 -- 2012}
%\institution[БГУИР]{Белорусский государственный университет информатики и радиоэлектроники}
%\logo{\includegraphics[width=1cm]{logo-kafEVM.eps}}

%
% Далее начинается сама презентация
%
\begin{document}
\begin{frame}
\titlepage
\begin{center}
e-mail: denis.pynkin@bsuir.by\\
{\tiny(если не отвечу на первое письмо -- не стесняйтесь прислать еще одно)}
\end{center}
\end{frame}

\section{Введение}

\begin{frame}{Цель дисциплины}
	Всестороннее изучение основных вопросов, связанных с функционированием сетевого программного обеспечения компьютерных сетей для различных архитектур и операционных систем.
\end{frame}

\begin{frame}{Задачи дисциплины}
	Подготовить специалиста в области сетевых технологий,  разбирающегося в принципах работы и умеющего создавать системное и прикладное сетевое программное обеспечение.
\end{frame}

\begin{frame}{Надо будет знать}
	\begin{itemize}
		\item основные возможности сетевых операционных систем
		\item основные протоколы обмена и интерфейсы, используемые при построении глобальных и корпоративных компьютерных сетей
		\item области применения,  достоинства и недостатки наиболее распространенных сетевых протоколов
		\item наиболее распространенные методы и алгоритмы взаимодействия программного обеспечения в компьютерных сетях
		\item принципы построения сетевого программного обеспечения
		\item особенности и принципы построения распределенных систем
	\end{itemize}
\end{frame}

\begin{frame}{Какого цвета учебник?}
	\begin{itemize}
		\item Стивенс У. Р.,  Феннер Б.,  Рудофф Э.М.,  UNIX: разработка сетевых приложений. 3-е изд. – СПб.:Питер,  2007,  1039 с.
		\item Тенненбаум Э.,  Ван Стеен М.,  Распределенные системы. Принципы и парадигмы,  Спб.:Питер,  2003,  877 с.
		\item Семенов Ю.А.,  Телекоммуникационные технологии [Электронный ресурс]. – Электронные данные. – Режим доступа: http://book.itep.ru/
	\end{itemize}
\end{frame}

\begin{frame}{Какого цвета учебник?}
	\begin{itemize}
		\item Тенненбаум Э.,  Компьютерные сети,  СПб.:Питер,  2003,  992 с.
		\item Под ред. Садыхова Р.Х.,  Средства параллельного программирования в ОС Linux: Учебное пособие,  Мн.: ЕГУ,  2004,  476 с.
		\item Камер Д.Э.,  Стивенс Д.Л.,  Сети TCP/IP,  том 3. Разработка приложений типа клиент/сервер для Linux/POSIX,  М.:Издательский дом "Вильямс",  2002,  592 с.
		\item Скляров И.С.,  Программирование боевого софта под Linux,  БХВ-Петербург,  2007,  416 с.
		\item Реймонд Эрик С.,  Искусство программирования для UNIX.: Пер. С англ. – М.: Издательский дом «Вильямс»,  2005,  544 с.
	\end{itemize}
\end{frame}



\begin{frame}{О чем курс?}
	\pause
	\begin{block}{}<2>
Уж точно не о сетевом администрировании!
	\end{block}
	\visible<3->
	\pause
	\begin{itemize}
		\item Стек протоколов TCP/IP
			\pause
		\item Сокеты -- сетевой API
			\pause
		\item Архитектура клиент-сервер
			\pause
		\item Введение в распределенные системы
			\pause
		\item MPI -- промышленный стандарт интерфейса для передачи сообщений
			\pause
		\item Всякая околосетевая фигня с точки зрения программиста
	\end{itemize}
\end{frame}

\begin{frame}{Примеры}
\begin{center}
	Все примеры рассчитаны, в первую очередь, на работу в среде POSIX
\par
	Проверялись на различных версиях ОС Linux, но должны работать на всех ОС, поддерживающих спецификацию POSIX (исключения будут оговариваться отдельно)
\par	
	В основном используется язык Си, как наиболее приближенный к системному уровню
\end{center}
\end{frame}

\begin{frame}{FAQ по работе примеров}
	\begin{itemize}
		\item Q: А вот почему в моей любимой {\footnotesize <здесь подставить название вашей любимой \sout{говно}ОС>} не работает ваш пример?\\
	A: Смотрите документацию по API к своей ОС!
		\item Q: А вот почему в моём любимом языке {\footnotesize <здесь подставить название вашего любимого языка>} не получается реализовать {\footnotesize <здесь подставить название фичи, которую не получается использовать>}?\\
	A1: Смотрите документацию к своему языку!\\
	A2: А вы точно уверены, что выбранный вами язык вообще может работать на этом уровне?\\
	A3: А вы точно уверены, что вы знаете выбранный вами язык программирования?\\
	A4: И вообще с чего вы взяли, что я знаю этот язык программирования?
	\end{itemize}
\end{frame}

\section{Введение в сетевое ПО}
\begin{frame}{Тест}
	тест
\end{frame}

\begin{frame}{}
\end{frame}

\begin{frame}{}
\end{frame}

\end{document}
%Конец файла
